% !TeX root = ../main.tex
% 系统需求分析

\chapter{系统需求分析}
    本章对多维度覆盖率融合显示系统的需求进行分析，主要包括功能需求与非功能需求。通过用户角色与使用场景分析，结合运行环境与约束条件，明确系统的目标与功能边界。

\section{系统总体需求}
    多维度覆盖率融合显示系统旨在为测试人员和质量保证人员提供统一的覆盖率分析平台。该系统部署于 SmartCI 平台，集成在冒烟测试流水线中。在冒烟测试执行后，采集多维度的代码覆盖率数据，进行细粒度的分析与可视化，从而实现以下目标：
    \begin{itemize}
        \item 支持测试人员在冒烟测试后快速查看代码覆盖率结果。
        \item 支持质量保证人员进行覆盖率的长期质量监控。
        \item 提供多维度的覆盖率展示，通过时间、领域、模块等分层次展示覆盖率数据，支持核心模块覆盖率和增量覆盖率展示。
    \end{itemize}

\section{功能需求}
\subsection{用例设计}
    多维度覆盖率融合显示系统由于目标硬件平台不同，需要支持用户在可视化界面配置自己的解析参数，主要有：
    \begin{itemize}
        \item \textbf{模块信息}：单板中每个领域的测试人员均需要配置模块信息，配置每个模块包含的文件以及该模块是否需要重点关注。这是细粒度覆盖率分析的重要依据。
        \item \textbf{排除文件}：覆盖率采集统计的是整个代码库中的代码覆盖情况，而特定领域只关注特定的代码文件，其余代码文件的覆盖率需要从覆盖率数据中排除，保持该领域代码覆盖率的独立性。
        \item \textbf{排除函数}：即使对于同一个代码文件，其中的部分函数对特定领域而言也不需要关注，例如维测函数或其他领域的函数，这些函数在执行测试用例时不会覆盖到，或不应在覆盖率统计中出现，需要排除。
        \item \textbf{二进制文件路径}：在 LLVM 覆盖率解析流程中，覆盖率的准确提取依赖于交叉编译生成的可执行文件、动态链接库以及可重定位对象文件。LLVM 的覆盖率解析工具在解析 \verb|.profraw| 文件时，需要结合对应的二进制文件，映射执行轨迹到具体的代码行，获取精确的覆盖率信息。
    \end{itemize}

    除了配置解析参数，系统还需要支持用户上传覆盖率数据包与构建包来解析覆盖率数据。其中，对于 LLVM 工具链，解析依赖的二进制文件保存在交叉编译得到的构建包中；而对于 GCC 工具链，交叉编译得到的源代码映射信息包含在 \verb|.gcno| 文件中，该文件会被自动包含在覆盖率采集模块导出的覆盖率数据包中，此时用户无需额外上传原始构建包。

    在解析完成后，用户能够查看多个层级的覆盖率数据，包括不同单板的融合覆盖率，单板级领域覆盖率，领域级模块覆盖率。其中每个层级均支持按时间维度展示，即同一维度下可呈现多个时间点的覆盖率柱状图，实现跨时间比较与趋势分析。此外，除了上述全量覆盖率视图，系统还提供增量覆盖率、核心模块覆盖率以及核心模块增量覆盖率等多维度展示，帮助用户全面掌握测试覆盖情况。

    系统应汇总并呈现覆盖率告警信息，辅助用户快速定位潜在问题，包括：文件未匹配模块告警、文件重复匹配模块告警、函数未覆盖告警。
    
    测试人员视角的总体用例视图如图~\ref{fig:3.1}所示。
    \begin{figure}[ht]
        \centering
        \includegraphics[width=0.8\textwidth]{usecases/测试人员-总体.pdf}
        \caption{测试人员-总体用例图}
        \label{fig:3.1}
    \end{figure}

    该用例图展示了整个系统的总体用例情况，主要区分为两个模块：POST操作与GET操作，即上传数据与获取数据。其中POST包括配置解析参数、上传覆盖率数据，GET包括查看告警信息、查看覆盖率数据。

\subsubsection{配置解析参数}
    配置解析参数的用例视图如图~\ref{fig:3.2}所示。
    \begin{figure}[ht]
        \centering
        \includegraphics[width=0.8\textwidth]{usecases/测试人员-配置解析参数.pdf}
        \caption{测试人员-配置解析参数用例图}
        \label{fig:3.2}
    \end{figure}

    该用例图展示了“配置解析参数”时的详细用例，包括配置模块信息，配置排除文件，配置排除函数，配置二进制文件，配置基准节点，配置管理员ID：
    \begin{itemize}
        \item \textbf{配置排除文件}：用户可指定在覆盖率分析中需要排除的文件，以保证特定领域覆盖率统计仅关注该领域相关代码。
        \item \textbf{配置排除函数}：用户可指定需要排除的函数，从而在覆盖率统计中忽略与当前领域无关或不需关注的函数，实现更精准的覆盖率分析。
        \item \textbf{配置模块信息}：本系统将一个单板的功能组件记为领域，而一个领域中又可以划分出许多功能模块。因此配置模块信息时应以领域为单位，模块信息包括领域名，领域责任人，以及领域包含的模块。其中，每个模块的信息包括模块责任人，模块是否为核心模块，以及模块包含的文件。
        \item \textbf{配置二进制文件（可选）}：在使用 LLVM 工具链时，覆盖率数据文件 \verb|.profraw| 不包含源代码行映射信息，需要用户提供交叉编译生成的二进制文件，其通常以压缩包形式存放在构建包中。用户需要配置压缩包路径和解压后的二进制文件路径，以便系统正确解析覆盖率数据到源代码行。
        \item \textbf{配置基准节点}：用于指定增量覆盖率分析的参考点，系统将以此为基准，计算最新节点相对于基准节点的增量覆盖率。在实际应用中，该基准节点可能是版本构建时的提交节点，也可能是用户自行指定的提交节点，支持不同场景下的覆盖率分析。
        \item \textbf{配置管理员ID}：用于指定有权限执行覆盖率解析与参数配置操作的管理员。每个领域可指定不同的管理员，以确保只有授权用户才能修改该领域的解析参数，避免误操作影响其他领域的配置。
    \end{itemize}

    注意到该用例图中多处标注了“可选”，其原因在于不同单板、不同领域的覆盖率解析行为不同：
    \begin{itemize}
        \item 某些单板或领域使用 GCC 工具链进行覆盖率采集，与 LLVM 工具链不同，其源代码行映射信息包含在 \verb|.gcno| 文件中，该文件会被覆盖率采集模块自动整合到覆盖率数据包中，无需上传构建包来解析二进制文件。此时二进制文件标注为可选。
        \item 某些领域较为特殊，虽然使用了 LLVM 工具链，但不存在完整的动态链接库等文件，而是采用编译过程中的某些中间文件实现了源代码行的映射，这些文件不会压缩在 IMG 中，不需要配置压缩包路径，只需要配置二进制文件路径即可。此时压缩包路径标注为可选。
    \end{itemize}

\subsubsection{查看告警信息}
    查看告警信息的用例视图如图~\ref{fig:3.3}所示。

    该用例图展示了查看告警信息的详细用例，告警信息包括模块信息告警和覆盖率告警，其中模块信息告警是为了警告用户模块信息配置出现异常，而覆盖率告警则是将未覆盖的信息聚合显示，便于测试人员定位问题。
    \begin{itemize}
        \item \textbf{模块告警}：包括文件未匹配模块告警和文件重复匹配模块告警；文件未匹配模块会导致某些文件未被模块统计，文件重复匹配模块会导致某些文件的覆盖率被记录了多次，二者均会导致模块覆盖率不准确。
        \item \textbf{覆盖率告警}：重点关注函数未覆盖告警，即覆盖率为0的函数。函数是软件系统的基本功能单元，函数级别的未覆盖意味着对应功能完全未验证，风险最大；相比之下，行或分支级别的覆盖缺失往往仅表示部分路径未执行，带来的风险相对较低。
    \end{itemize}

    \begin{figure}[ht]
        \centering
        \includegraphics[width=0.8\textwidth]{usecases/测试人员-查看告警信息.pdf}
        \caption{测试人员-查看告警信息用例图}
        \label{fig:3.3}
    \end{figure}

\subsubsection{查看覆盖率数据}
    查看覆盖率数据的用例视图如图~\ref{fig:3.4}所示。

    该用例图在测试人员视角与质量保证人员视角展示了查看覆盖率数据的详细用例。系统的覆盖率数据分为两个层级，一层是全量/增量覆盖率，第二层是单板分覆盖率、领域全模块覆盖率、领域核心模块覆盖率、覆盖率报告。
    
    增量覆盖率指的是当前代码节点（git commit id）与之前某个代码节点之间进行比较，期间增加的代码的覆盖率。与之对应的，全量覆盖率指的就是当前代码节点中所有代码的覆盖率。
    
    第二层级将覆盖率数据分为多个维度，每个维度粒度不同。系统主要通过柱状图展示覆盖率数据。
    \begin{itemize}
        \item \textbf{单板分领域覆盖率}：对于特定单板，统计其内部各个领域的整体覆盖率。此时柱状图的横坐标表示领域，纵坐标表示该领域的综合覆盖率。
        \item \textbf{领域全模块覆盖率}：对于特定领域，统计其内部所有模块的覆盖率。此时柱状图的横坐标表示模块，纵坐标表示该模块的覆盖率。
        \item \textbf{领域核心模块覆盖率}：根据模块信息中配置的核心模块标记，提取出需要重点关注的核心模块进行独立统计。此时，柱状图的含义随展示场景而变化：
            \begin{itemize}
                \item 在单板维度，柱状图纵坐标表示某领域内所有核心模块的聚合覆盖率；
                \item 在领域维度，柱状图纵坐标表示单个核心模块的覆盖率。
            \end{itemize}
        通过这种区分，系统能够突出关键模块的质量状况，辅助用户优先优化关键功能模块。
        \item \textbf{覆盖率报告}：所有柱状图均支持下钻操作：用户点击柱子后，可跳转至相应的覆盖率报告。报告由 LCOV 工具生成的 HTML 页面呈现，能够在代码级别直观标注每个文件的覆盖情况（是否覆盖，覆盖次数等），为用户提供最细粒度的覆盖率证据。
    \end{itemize}

    \begin{figure}[ht]
        \centering
        \includegraphics[width=0.8\textwidth]{usecases/测试-质量-查看覆盖率数据.pdf}
        \caption{测试-质量-查看覆盖率数据用例图}
        \label{fig:3.4}
    \end{figure}

    该用例图中，使用自定义标签\verb|<<time-related>>|标注了部分用例，表示这些用例支持获取不同时间点的覆盖率数据，并聚合对比展示，实现跨时间比较与趋势分析。

\subsection{需求分析}
    功能需求分析主要关注系统应当具备的功能特性，系统主要功能点包括 覆盖率数据处理，覆盖率可视化呈现，自动化测试流程。

\subsubsection{覆盖率数据处理}
    系统能够从不同硬件平台（单板）采集覆盖率数据，实现多维度的覆盖率分析与融合，并实现覆盖率数据与覆盖率报告的持久化存储。
    \begin{itemize}
        \item \textbf{功能描述}：
        \begin{itemize}
            \item \textbf{解析参数配置}：允许用户自定义覆盖率解析参数，支持多维度数据解析；通过界面权限控制，确保解析配置不被误操作修改。
            \item \textbf{覆盖率采集}：兼容多种交叉编译工具链（LLVM、GCC），从冒烟测试流程中获取插桩统计的覆盖率数据，并支持多单板数据统一处理。
            \item \textbf{细粒度分析}：解析不同工具链采集的原始覆盖率数据，按单板、领域、模块分层统计，支持增量覆盖率计算及多时间点对比和趋势分析。
            \item \textbf{持久化存储}：将解析得到的覆盖率数据保存至数据库，将生成的 HTML 覆盖率报告存储到静态文件系统。
        \end{itemize}
        \item \textbf{用户角色}：测试人员。
        \item \textbf{输入/输出}：
        \begin{itemize}
            \item 输入：冒烟测试流程产生的插桩覆盖率数据。
            \item 输出：持久化存储的覆盖率数据及覆盖率报告。
        \end{itemize}
    \end{itemize}


\subsubsection{覆盖率可视化展示}
    系统提供覆盖率数据可视化功能，以帮助测试人员和质量保证人员快速理解覆盖率情况并发现潜在问题。。
    \begin{itemize}
        \item \textbf{功能描述}：
            \begin{itemize}
                \item \textbf{展示多维度覆盖率数据}：展示不同维度的全量覆盖率与增量覆盖率，并提供报告查看功能，支持按时间、单板、领域、模块等维度进行筛选与对比。
                \item \textbf{展示告警信息}：展示模块配置告警与函数未覆盖告警，帮助用户发现潜在问题。
            \end{itemize}
        \item \textbf{用户角色}：测试人员，质量保证人员。

        \item \textbf{输入/输出}：
        \begin{itemize}
            \item 输入：解析得到的覆盖率数据及覆盖率报告。
            \item 输出：多维度覆盖率可视化界面，包括图表展示和告警列表等。
        \end{itemize}
    \end{itemize}

\subsubsection{自动化测试流程}
    系统集成在 SmartCI 持续集成平台的冒烟测试流水线中，实现覆盖率数据的自动化采集分析与可视化，确保高效性，减少人工干预。主要功能如下：
    \begin{itemize}
        \item \textbf{功能描述}：
        \begin{itemize}
            \item \textbf{自动化执行}：系统在冒烟测试流程中，自动采集覆盖率数据。该过程通过 SmartCI 流水线任务调度，实现定时或手动触发，支持多节点并行执行。
            \item \textbf{结果回传与处理}：测试完成后，系统自动将覆盖率数据回传至主机，经过解析生成标准化覆盖率信息，并存储于数据库，同时生成 HTML 覆盖率报告，供后续可视化分析使用。
        \end{itemize}
        \item \textbf{用户角色}：测试人员。
        \item \textbf{输入/输出}：
        \begin{itemize}
            \item 输入：冒烟测试流程产生的插桩覆盖率数据。
            \item 输出：持久化存储的覆盖率数据及覆盖率报告。
        \end{itemize}
    \end{itemize}

\section{非功能需求分析}
    Roman\cite{Roman1985ATO} 指出非功能需求分析在软件开发中具有保障系统质量和可持续发展的重要作用。系统除了满足功能性需求外，还需满足以下非功能需求：
    \begin{itemize}
        \item \textbf{性能}：系统应保证在常规数据规模下，覆盖率数据的细粒度分析以及可视化呈现延迟不超过合理阈值。
        \item \textbf{可用性}：前端界面应提供清晰、易操作的覆盖率数据浏览和分析功能，保证用户能够快速获取所需信息。
        \item \textbf{可扩展性}：系统设计应支持未来增加更多硬件平台、领域或工具链的覆盖率数据处理能力。
        \item \textbf{安全性}：对覆盖率数据和生成的报告提供合理的访问控制和数据保护措施，确保敏感信息不被未授权访问。
        \item \textbf{可维护性}：系统结构应清晰，代码规范，便于后续功能扩展和维护。
    \end{itemize}

\section{本章小结}
    本章对多维度覆盖率融合显示系统进行了系统需求分析。功能需求分析中，通过用例设计明确了系统主要用户及其交互行为，并对核心功能进行了需求分析。非功能需求分析中，对性能、可靠性、可扩展性和安全性提出了约束。总体上，本章为后续概要设计阶段提供了明确的需求基础。
